\documentclass[12pt,fleqn]{article}

% --- Настройка шрифтов для XeLaTeX ---
\usepackage{fontspec} 
\usepackage[russian]{babel}

% В XeLaTeX кодировка UTF-8 и поиск символов работают нативно.
% Пакеты inputenc и cmap здесь НЕ НУЖНЫ и вызывают ошибки.

\setmainfont{Liberation Serif}
\setsansfont{Liberation Sans}
\setmonofont{Liberation Mono}

% --- Стандартные пакеты ---
\usepackage[a4paper, left=1cm, right=1cm, top=1.27cm, bottom=1.27cm]{geometry}
\usepackage{graphicx}
\usepackage{array}
\usepackage{setspace}
\usepackage{tabularx}
\usepackage{makecell}
\usepackage{longtable}
\usepackage{multirow}
\usepackage{caption}
\usepackage{amsmath}
\usepackage{amssymb}
\usepackage{ragged2e}

\onehalfspacing


\begin{document}

\begin{flushleft}
    Домашнее задание №1 \\
    Пахомов Владимир Юрьевич 467036
\end{flushleft}

\begin{justify}
    \section{Условие (Вариант 22)}
    \begin{itemize}
    \item Схема 2
    \item $\uparrow J_5$ = 1.3 А
    \item $\leftarrow E_2$ = 32 В
    \item $\downarrow E_4$ = 11,5 В
    \item $R_1 $ = 8 Ом
    \item $R_2 $ = 9 Ом
    \item $R_3 $ = 8 Ом
    \item $R_4 $ = 5 Ом
    \end{itemize}
    а) скомпоновать схему согласно своему варианту; \\
    б) найти все неизвестные токи, используя законы Кирхгофа; \\
    в) найти все неизвестные токи, используя метод контурных токов; \\
    г) найти все неизвестные токи, используя метод узловых напряжений; \\
    д) Определить напряжение, приложенное к источнику тока.; \\
    е) осуществить проверку, составив баланс мощностейж \\
    ж) рассчитать ток любой ветви, содержащей источник ЭДС методом эквивалентных преобразований или методом эквивалентного генератора. \\

\newpage

    \section{Решение}
    \subsection{а) Компоновка схемы}
    \begin{center}
        \includegraphics[height=10cm]{schema1.png}
    \end{center}

    \subsection{б) Законы Кирхгофа}
    \begin{itemize}
        \item $p^* = 5 \;- $ количество ветвей
        \item $ p_\text{ит} = 1 \; - $ ветви с известными токами
        \item $ p = p* - p_\text{ит} = 4 \; - $ количество неизвестных токов
        \item $ q = 3 \; - $ количество узлов
        \item $ n = p - (q - 1) = 2 \; - $ количество независимых контуров
        \item $ m_I = q - 1 = 2 \; - $ количество уравнений для ЗКI
        \item $ m_{II} = n = 2 \; - $ количество уравнений для ЗКII
    \end{itemize}
    Составим уравнения для узлов 1 и 3 и для двух контуров:
    \begin{equation}
        \begin{cases}
            I_1 = I_2 + I_3 \\
            I_2 + I_4 = J \\
            E_2 = R_1I_3 + R_2I_1 \\
            E_2 - E_4 = R_2I_1 + R_3I_2 - R_4I_4 \\
        \end{cases}
    \end{equation}
    Подставив занчения сопротивлений и ЭДС получим:
    \begin{equation}
        \begin{cases}
            I_1 = I_2 + I_3 \\
            I_2 + I_4 = 1.3 \\
            32 = 8I_3 + 9I_1 \\
            32 - 11.5 = 9I_1 + 8I_2 - 5I_4 \\
        \end{cases}
    \end{equation}
    Решив эту систему получаем ответ:
    \begin{equation}
        \begin{cases}
            I_1 = 2.157 \; \text{А} \\
            I_2 = 0.584 \; \text{А} \\
            I_3 = 1.573 \; \text{А} \\
            I_4 = 0.716 \; \text{А} \\
        \end{cases}
    \end{equation}
    
\newpage

    \subsection{в) Метод контурных токов}

    \begin{itemize}
        \item $p^* = 5 \;- $ количество ветвей
        \item $ p_\text{ит} = 1 \; - $ ветви с известными токами
        \item $ p = p* - p_\text{ит} = 4 \; - $ количество неизвестных токов
        \item $ q = 3 \; - $ количество узлов
        \item $ n = p - (q - 1) = 2 \; - $ количество неизвестных контурных токов
        \item $ m = p_\text{ит} = 1 \; - $ количество известных контурных токов
    \end{itemize}

    $I_{33} = J = 1.3 \text{А}$ \\
    Составим систему уравнений:
    \begin{equation}
        \begin{cases}
            R_{11}I_{11} + R_{12}I_{22} + R_{13}I_{33} = E_{11} \\
            R_{21}I_{11} + R_{22}I_{22} + R_{23}I_{33} = E_{22} \\
        \end{cases}
    \end{equation}
    Посчитав ЭДС и сопротивления контуров получаем:
    \begin{equation}
        \begin{cases}
            (R_1 + R_2)I_{11} + R_2I_{22} + 0I_{33} = E_2 \\
            R_2I_{11} + (R_2 + R_3 + R_4)I_{22} - R_4I_{33} = E_2 = E_4 \\
        \end{cases}
    \end{equation}
    \begin{equation}
        \begin{cases}
            17I_{11} + 9I_{22} = 32 \\
            9I_{11} + 22I_{22} - 5 \cdot 1.3 = 32 - 11.5 \\
        \end{cases}
    \end{equation}

    Решив систему получим: \\
    $ I_{11} \approx 1.573 \text{А} $ \\
    $ I_{22} \approx 0.584 \text{А} $ \\

    Тогда: \\
    $ I_1 = I_{11} + I_{22} = 1.573 + 0.584 = 2.157 \text{А} $ \\
    $ I_2 = I_{22} = 0.584 \text{А} $ \\
    $ I_3 = I_{11} = 1.573 \text{А} $ \\
    $ I_4 = I_{33} - I_{22} = 1.3 - 0.584 = 0.716 \text{А} $ \\

    Ответ: \\
    \begin{equation}
        \begin{cases}
            I_1 = 2.157 \; \text{А} \\
            I_2 = 0.584 \; \text{А} \\
            I_3 = 1.573 \; \text{А} \\
            I_4 = 0.716 \; \text{А} \\
        \end{cases}
    \end{equation}

\newpage

\subsection{г) Метод узловых напряжений}
    \begin{itemize}
        \item $p^* = 5 \;- $ количество ветвей
        \item $ p_\text{ит} = 1 \; - $ ветви с известными токами
        \item $ p = p* - p_\text{ит} = 4 \; - $ количество неизвестных токов
        \item $ q = 3 \; - $ количество узлов
        \item $ l = q - 1 = 2 \; -$ количество узловых напряжений 
    \end{itemize}

    Произвольно обозначим узловые напряжения $U_{10}$ и $U_{20}$.
    Примем узел 3 за заземленный, тогда его потенциал $\varphi_3 = 0$.
    Узловые напряжения направлены от незаземленных узлов к заземленному.\\
    Составим систему уравнений: \\
    \begin{equation}
        \begin{cases}
            g_{11}U_{10} + g_{12}U_{20} = J_{11} \\
            -g_{21}U_{10} + g_{22}U_{20} = J_{22} \\
        \end{cases}
    \end{equation}

    \begin{itemize}
        \item $ g_{11} = \frac{1}{R_1} + \frac{1}{R_2} + \frac{1}{R_3} \approx 0.3611 $
        \item $ g_{22} = \frac{1}{R_1} + \frac{1}{R_2} + \frac{1}{R_4} \approx 0.4361 $
        \item $ g_{12} = g_{21}= \frac{1}{R_1} + \frac{1}{R_2} \approx 0.2361 $
        \item $ J_{11} = \frac{E_2}{R_2} = \approx 3.5556 \; \text{А}$
        \item $ J_{22} = J - \frac{E_2}{R_2} = \approx -4.5556 \; \text{А}$
    \end{itemize}

    Подставив в систему получим: \\
    \begin{equation}
        \begin{cases}
            0.3611U_{10} - 0.2361U_{20} = 3.5556 \\
            -0.2361U_{10} + 0.4361U_{20} = -4.5556 \\        
        \end{cases}
    \end{equation}

    После решения системы получаем значения $U_{10} \approx 4.669 \; \text{В} $ и $U_{20} \approx -7.919 \; \text{В} $.\\

    Для определения токов воспользуемся законом ома $I = \frac{\sum \pm U + \sum \pm E}{\sum R}$:
    \begin{equation}
        \begin{cases}
            I_1 = \frac{U_{20} - U_{10} - E_2}{R_2} \approx 2.157 \; \text{А} \\
            I_2 = \frac{U_{10}}{R_3} \approx 0.584 \; \text{А} \\
            I_3 = \frac{U_{10} - U_{20}}{R_1} \approx 1.574 \; \text{А} \\
            I_4 = \frac{U_{20} + E_4}{R_4} \approx 0.716 \; \text{А} \\
        \end{cases}
    \end{equation}
    

    \newpage

    \subsection{д) Определение напряжения, приложенного к источнику тока}
    Возьмем третий (голубой) контур с источником тока $J$. Уравнение будет иметь ввид:
    \[-U_J - I_4 \cdot R_4 - E_4 = 0 \Rightarrow U_J = -E_4 - I_4 \cdot R_4 \Rightarrow U_J = -15.08 \; \text{В} \]
    Отрицательный знак означает, что потенциал нижнего узла (3) ниже потенциала верхнего узла (2) и направление напряжения противоположно выбранному.\\

    \subsection{е) Определение мощностей источников}
    Мощность источника ЭДС положительна, если направление тока через него совпадает
    с направлением ЭДС. Мощность источника тока определяется через напряжение на его зажимах.

    \begin{itemize}
        \item $ P_{E_2} = E_2 \cdot I_1 = 69.024 \; \text{Вт} $
        \item $ P_{E_4} = E_4 \cdot I_4 = 8.234 \; \text{Вт} $
        \item $ P_J = U_J \cdot J = -10.296 \; \text{Вт} $
    \end{itemize}

    Суммарная мощность источников:
    \[ P_\text{и} = P_{E_2} + P_{E_4} + P_J = 66.962 \; \text{Вт} \]

    Мощность на резистивных элементах:
    \begin{itemize}
        \item $ P_{R_1} = R_1 \cdot I_3^2 = 19.816 \; \text{Вт} $
        \item $ P_{R_2} = R_2 \cdot I_1^2 = 41.877 \; \text{Вт} $
        \item $ P_{R_3} = R_3 \cdot I_2^2 = 2.728 \; \text{Вт} $
        \item $ P_{R_4} = R_4 \cdot I_4^2 = 2.565 \; \text{Вт} $
    \end{itemize}

    Суммарная мощность потребителей:
    \[ P_\text{п} = P_{R_1} + P_{R_2} + P_{R_3} + P_{R_4} = 66.986 \; \text{Вт} \]
    
    Условие равенства почти выполняется: $P_\text{и} = P_\text{п}$. Погрешность составляет
    менее $1\%$. Значит можно сказать, что баланс мощностей сошелся.

    \newpage

    \subsection{ж) Рассчет тока ветви с источником $E_2$ методом эквивалентного генератора}

    \begin{center}
        \includegraphics[height=8cm]{schema2.png}
    \end{center}

    Для определения $E_\text{г}$ разомкнем ветвь с искомым током $I_1$. ЭДС генератора равна
    напряжению холостого хода между узлами 2 и 1:

    Примем $\varphi_3 = 0$
    \begin{itemize}
        \item Узел 1: $ \frac{\varphi_3}{R_3} + \frac{\varphi_1 - \varphi_2}{R_1} = 0 \Rightarrow \varphi_2 = 2\varphi_1 $
        \item Узел 2: $ J = \frac{\varphi_2 - \varphi_1}{R_1} + \frac{\varphi_2 - (-E_4)}{R_4} \Rightarrow 1.3 = \frac{2\varphi_1 - \varphi_1}{8} + \frac{\varphi_2 - (-E_4)}{5} \Rightarrow \varphi_1 \approx -1.905 \; \text{В} \quad \varphi_2 \approx -3.81 \; \text{В} $
    \end{itemize}

    Напряжение холостого хода:
    \[ E_\text{г} = U_{21xx} = \varphi_2 - \varphi_1 = -1.905 \; \text{В} \]

    Для нахождения $R_\text{г}$ составим пассивную цепь. Источник ЭДС $E_4$ закорачиваем, источник тока $J$ обрываем.
    Резисторы $R_3$ и $R_4$ оказываются соединенными последовательно и они параллельны $R_1$. Тогда:
    \[ R_\text{г} = \frac{R_1(R_3 + R_4)}{R1 + R_3R_4} \approx 4.952 \; \text{Ом} \]

    Для нахождения искомого $I_1$ составим цепь, содержащую параемтры генератора $E_\text{г}$ и $R_\text{г}$ и параметры удаленной ветви $E_2$ и $R_2$.
    По второму закону Кирхгофа для контура: 
    \[ (R_2 + R_\text{г}) \cdot I_1 = E_2 + E_\text{{г}} \Rightarrow I_1 = \frac{E_2 + E_\text{г}}{R_2 + R_\text{г}} \approx 2.157 \; \text{А} \]
    Значение совпало с результатами других методов.

\end{justify}


\end{document}
